% ============================================================
%  Министерство образования и науки
%  [ITHUB]
%  Кафедра [Информационная безопасность]
%  Направление: [DevOps]
%  Лабораторная работа № [12]
%  Локальная разработка. Линейные алгоритмы
%  Сборка: pdflatex main.tex  (запустить 2 раза)
% ============================================================
\documentclass[12pt,a4paper]{article}

\usepackage[T2A]{fontenc}
\usepackage[utf8]{inputenc}
\usepackage[russian]{babel}
\usepackage{geometry}
\geometry{left=30mm,right=15mm,top=20mm,bottom=20mm}
\usepackage{setspace}
\onehalfspacing
\usepackage{amsmath}
\usepackage{array}
\usepackage{fancyvrb}
\usepackage{enumitem}

\fvset{fontsize=\small}
\setlength{\parindent}{1.25cm}
\setlength{\parskip}{0pt}

\begin{document}

% ------------------------------------------------------------
% 1. ТИТУЛЬНЫЙ ЛИСТ
% ------------------------------------------------------------
\begin{titlepage}
\begin{center}
Министерство образования и науки\\
\textbf{[ITHUB]}\\
Кафедра [Информационная безопасность]\\
Направление: [DevOps]

\vspace{2.5cm}
{\Large\textbf{Лабораторная работа № [12]}}\$$0.5cm]
{\Large\textbf{Локальная разработка. Линейные алгоритмы}}

\vspace{1.5cm}
{\large Тема: разработка программы на языке C++ для решения\\
задачи 2.25 из сборника задач А.\,Г.\,Юркина}\$$0.3cm]
{\large\textit{<<На шахматной доске расположены белый король,\\
чёрная ладья и слон>>}}
\end{center}

\vspace{2cm}
\begin{tabular}{p{5cm}p{7cm}}
Выполнил(а), группа: & \underline{\hspace{7cm}} \$$0.4cm]
Проверил(а):         & \underline{\hspace{7cm}} \$$0.4cm]
Оценка:              & \underline{\hspace{7cm}} \$$0.4cm]
Дата защиты:         & \underline{\hspace{7cm}} \\
\end{tabular}

\vfill
\begin{center}
2026 г.
\end{center}
\end{titlepage}

% ------------------------------------------------------------
% 2. ПОСТАНОВКА ЗАДАЧИ
% ------------------------------------------------------------
\section*{2. Постановка задачи}

\textbf{Задача 2.25} из сборника задач А.\,Г.\,Юркина (раздел~2):

\begin{quote}
\textit{<<На шахматной доске расположены белый король, чёрная ладья
и слон. Определить, есть ли угроза королю, если да, то от кого?
Слон ходит по диагонали. Ладья --- по горизонталям и вертикалям.>>}
\end{quote}

\textbf{Дано:} шахматная доска $8\times 8$; расположение трёх фигур:
белого короля, чёрной ладьи и чёрного слона.

\textbf{Требуется:} определить, находится ли белый король под боем
хотя бы одной из чёрных фигур; если да --- указать, какая фигура
угрожает королю (ладья, слон или обе).

\textbf{Формат ввода:} координаты трёх фигур --- три пары целых чисел
$(x, y)$, где $x$ --- номер вертикали ($a=1,\ldots,h=8$), $y$ --- номер
горизонтали ($1,\ldots,8$).

\textbf{Формат вывода:} одно из четырёх сообщений:
\begin{itemize}[nosep]
  \item <<Угроза королю от ладьи>>;
  \item <<Угроза королю от слона>>;
  \item <<Угроза королю от ладьи и слона>>;
  \item <<Угрозы королю нет>>.
\end{itemize}

% ------------------------------------------------------------
% 3. ВЕРБАЛЬНАЯ МОДЕЛЬ РЕШЕНИЯ
% ------------------------------------------------------------
\section*{3. Вербальная модель решения}

Словесное описание алгоритма:

\begin{enumerate}
  \item Ввести координаты трёх фигур: белого короля ($K$), чёрной
        ладьи ($R$) и чёрного слона ($B$). Проверить, что все
        координаты лежат в пределах доски (от 1 до 8) и что фигуры
        не стоят на одной клетке. При нарушении --- сообщение об
        ошибке и завершение работы.
  \item Проверить, угрожает ли ладья королю: ладья бьёт по вертикали
        и горизонтали, поэтому угроза возможна только если король и
        ладья стоят на одной вертикали (равны абсциссы) или на одной
        горизонтали (равны ординаты).
  \item Если общая линия есть, проверить, не перекрывает ли её слон:
        если слон стоит строго между королём и ладьёй на этой же
        линии, угрозы от ладьи нет; иначе зафиксировать угрозу от ладьи.
  \item Проверить, угрожает ли слон королю: слон бьёт по диагонали,
        поэтому угроза возможна только если модуль разности абсцисс
        короля и слона равен модулю разности их ординат (фигуры не
        совпадают).
  \item Если общая диагональ есть, проверить, не перекрывает ли её
        ладья: если ладья стоит строго между королём и слоном на этой
        же диагонали, угрозы от слона нет; иначе зафиксировать угрозу
        от слона.
  \item По двум признакам вывести заключение:
        \begin{itemize}[nosep]
          \item угрожают обе фигуры --- <<Угроза от ладьи и слона>>;
          \item угрожает только ладья --- <<Угроза от ладьи>>;
          \item угрожает только слон --- <<Угроза от слона>>;
          \item не угрожает никто --- <<Угрозы нет>>.
        \end{itemize}
  \item Завершить работу программы.
\end{enumerate}

\textbf{Примечание.} Правила перекрытия (шаги 3 и 5) формализуют
шахматное правило: фигура не может бить сквозь другую фигуру. Если
требуется упрощённая постановка (перекрытием пренебречь), шаги 3 и~5
исключаются.

% ------------------------------------------------------------
% 4. МАТЕМАТИЧЕСКАЯ МОДЕЛЬ РЕШЕНИЯ
% ------------------------------------------------------------
\section*{4. Математическая модель решения}

\textbf{4.1. Обозначения.} Доску отождествим с множеством клеток
$$
S = \{\, (x,\,y) \mid x,\,y \in \{1, 2, \ldots, 8\} \,\},
$$
где $x$ --- вертикаль ($a=1,\ldots,h=8$), $y$ --- горизонталь.
Положение фигур:
$$
K = (x_K, y_K), \qquad R = (x_R, y_R), \qquad B = (x_B, y_B),
$$
причём $K \neq R$, $K \neq B$, $R \neq B$.

\textbf{4.2. Общая линия <<король -- ладья>> (условие боя ладьи):}
$$
L_R = \bigl( x_K = x_R \bigr) \lor \bigl( y_K = y_R \bigr).
$$

\textbf{4.3. Общая диагональ <<король -- слон>> (условие боя слона):}
$$
L_B = \bigl( |x_K - x_B| = |y_K - y_B| \bigr) \land
      \bigl( x_K \neq x_B \bigr).
$$

\textbf{4.4. Перекрытие линии атаки.} Фигура-перекрыватель должна
лежать строго между атакующей фигурой и королём на той же линии:
\begin{align*}
Z_R &= \bigl( x_B = x_K = x_R \ \land \
        \min(y_K,y_R) < y_B < \max(y_K,y_R) \bigr) \\
    &\quad \lor \bigl( y_B = y_K = y_R \ \land \
        \min(x_K,x_R) < x_B < \max(x_K,x_R) \bigr), \$$4pt]
Z_B &= \bigl( |x_R - x_K| = |y_R - y_K| \bigr)
      \land \bigl( |x_R - x_B| = |y_R - y_B| \bigr) \\
    &\quad \land \ \min(x_K,x_B) < x_R < \max(x_K,x_B)
      \ \land \ \min(y_K,y_B) < y_R < \max(y_K,y_B).
\end{align*}

\textbf{4.5. Признаки угрозы:}
$$
A_R = L_R \land \lnot Z_R, \qquad
A_B = L_B \land \lnot Z_B.
$$

\textbf{4.6. Итоговая функция.} Результат определяется вектором
$(A_R, A_B)$:

\begin{center}
\begin{tabular}{|c|c|l|}
\hline
$A_R$ & $A_B$ & Вывод программы \\
\hline
1 & 1 & Угроза королю от ладьи и слона \\
\hline
1 & 0 & Угроза королю от ладьи \\
\hline
0 & 1 & Угроза королю от слона \\
\hline
0 & 0 & Угрозы королю нет \\
\hline
\end{tabular}
\end{center}

\textbf{4.7. Оценка сложности.} Алгоритм линейный (без циклов):
число операций фиксировано и не зависит от входных данных;
сложность $O(1)$ по времени и $O(1)$ по памяти.

% ------------------------------------------------------------
% 5. БЛОК-СХЕМА
% ------------------------------------------------------------
\section*{5. Блок-схема (декомпозиция задачи)}

\begin{Verbatim}[frame=single]
                    +------------------------------+
                    |            START             |
                    +--------------+---------------+
                                   v
                    +------------------------------+
                    | Vvod K=(xK,yK), R=(xR,yR),   |
                    | B=(xB,yB)                    |
                    +--------------+---------------+
                                   v
                    +------------------------------+   net
                    | Koordinaty 1..8 i kletki     +-------+
                    | ne sovpadayut?               |       |
                    +--------------+---------------+       v
                                da v                 "Oshibka vvoda"
                    +------------------------------+       |
                    | LR := (xK=xR) OR (yK=yR)     |       |
                    +--------------+---------------+       |
                                   v                       |
                    +------------------------------+  da   |
                    | Slon mezhdu K i R            +------+|
                    | na etoy linii?               |      ||
                    +--------------+---------------+      ||
                                net v                      ||
                    +------------------------------+      ||
                    | AR := LR                     |<-----+|
                    +--------------+---------------+       |
                                   v                       |
                    +------------------------------+       |
                    | LB := |xK-xB| = |yK-yB|      |       |
                    +--------------+---------------+       |
                                   v                       |
                    +------------------------------+  da   |
                    | Ladya mezhdu K i B           +------+|
                    | na diagonali?                |      ||
                    +--------------+---------------+      ||
                                net v                      ||
                    +------------------------------+      ||
                    | AB := LB                     |<-----+|
                    +--------------+---------------+       |
                                   v                       |
                    +------------------------------+       |
                    | AR=1 i AB=1 ?                +-- da -> "Ugroza ot ladyi i slona"
                    +--------------+---------------+
                                   | net
                                   v
                    +------------------------------+
                    | AR=1 ?                       +-- da --> "Ugroza ot ladyi"
                    +--------------+---------------+
                                   | net
                                   v
                    +------------------------------+
                    | AB=1 ?                       +-- da --> "Ugroza ot slona"
                    +--------------+---------------+
                                   | net
                                   v
                          "Ugrozy korolyu net"
                                   v
                    +------------------------------+
                    |            STOP              |
                    +------------------------------+
\end{Verbatim}

% ------------------------------------------------------------
% 6. ПРОГРАММА (КЛЮЧЕВЫЕ МОМЕНТЫ)
% ------------------------------------------------------------
\section*{6. Программа на языке высокого уровня (ключевые моменты)}

Программа написана на языке C++ (стандарт C++11 и выше). Полный
листинг --- в приложении~А. Ниже разобраны ключевые фрагменты.

\textbf{6.1. Представление данных.} Положение фигуры --- структура
\texttt{Cell} из двух целых полей:

\begin{Verbatim}
struct Cell {
    int x;
    int y;
};
\end{Verbatim}

\textbf{6.2. Геометрические предикаты.} Условия $L_B$, $Z_R$, $Z_B$
математической модели реализованы функциями \texttt{onDiagonal},
\texttt{betweenStraight}, \texttt{betweenDiagonal}:

\begin{Verbatim}
bool onDiagonal(const Cell& a, const Cell& b) {
    return abs(a.x - b.x) == abs(a.y - b.y) && a.x != b.x;
}

bool betweenDiagonal(const Cell& p, const Cell& a, const Cell& b) {
    return onDiagonal(a, p) && onDiagonal(p, b) &&
           abs(a.x - p.x) < abs(a.x - b.x) &&
           abs(a.y - p.y) < abs(a.y - b.y);
}
\end{Verbatim}

\textbf{6.3. Признаки угрозы.} Каждая функция возвращает \texttt{bool}
и дословно повторяет формулу признака:

\begin{Verbatim}
bool rookThreat(const Cell& k, const Cell& r, const Cell& b) {
    if (k.x != r.x && k.y != r.y) return false;
    if (betweenStraight(b, k, r)) return false;
    return true;
}

bool bishopThreat(const Cell& k, const Cell& b, const Cell& r) {
    if (!onDiagonal(k, b))        return false;
    if (betweenDiagonal(r, k, b)) return false;
    return true;
}
\end{Verbatim}

\textbf{6.4. Вывод результата.} Ветвление по таблице п.\,4.6:

\begin{Verbatim}
bool rook   = rookThreat(K, R, B);
bool bishop = bishopThreat(K, B, R);

if (rook && bishop)  cout << "Угроза королю от ладьи и слона.\n";
else if (rook)       cout << "Угроза королю от ладьи.\n";
else if (bishop)     cout << "Угроза королю от слона.\n";
else                 cout << "Угрозы королю нет.\n";
\end{Verbatim}

\textbf{6.5. Контроль ввода.} Функция \texttt{readCell} проверяет
попадание координат в диапазон $1\ldots 8$; дополнительно проверяется,
что фигуры не стоят на одной клетке.

% ------------------------------------------------------------
% 7. ПРОВЕРКА
% ------------------------------------------------------------
\section*{7. Проверка}

Проверка выполнена на основе математической модели (п.\,4): для каждого
набора координат значения $A_R$ и $A_B$ вычислялись вручную по формулам
4.2--4.5 и сопоставлялись с выводом программы. Обозначения: $+$ ---
признак истинен, $-$ --- ложен, --- --- не определялся.

\begin{center}
\small
\begin{tabular}{|c|c|c|c|c|c|c|c|c|c|p{4.2cm}|}
\hline
№ & $K$ & $R$ & $B$ & $L_R$ & $Z_R$ & $A_R$ & $L_B$ & $Z_B$ & $A_B$ & Ожидаемый вывод \\
\hline
1 & (5,4) e4 & (5,8) e8 & (1,1) a1 & $+$ & $-$ & $+$ & $-$ & --- & $-$ & Угроза от ладьи \\
\hline
2 & (5,4) e4 & (1,1) a1 & (8,1) h1 & $-$ & --- & $-$ & $+$ & $-$ & $+$ & Угроза от слона \\
\hline
3 & (5,4) e4 & (1,1) a1 & (8,2) h2 & $-$ & --- & $-$ & $-$ & --- & $-$ & Угрозы нет \\
\hline
4 & (5,4) e4 & (5,8) e8 & (5,6) e6 & $+$ & $+$ & $-$ & $-$ & --- & $-$ & Угрозы нет (линию перекрыл слон) \\
\hline
5 & (5,4) e4 & (3,6) c6 & (1,8) a8 & $-$ & --- & $-$ & $+$ & $+$ & $-$ & Угрозы нет (диагональ перекрыта ладьёй) \\
\hline
6 & (5,4) e4 & (5,8) e8 & (1,8) a8 & $+$ & $-$ & $+$ & $+$ & $-$ & $+$ & Угроза от ладьи и слона \\
\hline
\end{tabular}
\end{center}

\textbf{Разбор контрольных примеров (по формулам п.\,4):}
\begin{itemize}
  \item \textbf{№\,1.} $x_K = x_R = 5 \Rightarrow L_R = 1$; слон
        $(1,1)$ не лежит на вертикали $x=5 \Rightarrow Z_R = 0
        \Rightarrow A_R = 1$. Для слона: $|5-1| = 4 \neq |4-1| = 3
        \Rightarrow L_B = 0 \Rightarrow A_B = 0$. Вывод: <<Угроза
        от ладьи>>.
  \item \textbf{№\,2.} $|5-8| = |4-1| = 3 \Rightarrow L_B = 1$;
        ладья $(1,1)$ не лежит на диагонали $K$--$B$, так как
        $|1-5| = 4 \neq |1-4| = 3 \Rightarrow Z_B = 0
        \Rightarrow A_B = 1$. Для ладьи: $5 \neq 1$, $4 \neq 1
        \Rightarrow L_R = 0$. Вывод: <<Угроза от слона>>.
  \item \textbf{№\,4.} $L_R = 1$, но слон $(5,6)$: $x_B = x_K = x_R
        = 5$ и $4 < 6 < 8 \Rightarrow Z_R = 1 \Rightarrow A_R = 0$;
        общей диагонали с королём нет ($x_K = x_B$). Вывод:
        <<Угрозы нет>>.
  \item \textbf{№\,5.} $|5-1| = |4-8| = 4 \Rightarrow L_B = 1$;
        ладья $(3,6)$: $|3-5| = |6-4| = 2$ --- на диагонали, $2 < 4
        \Rightarrow Z_B = 1 \Rightarrow A_B = 0$; $5 \neq 3$,
        $4 \neq 6 \Rightarrow L_R = 0$. Вывод: <<Угрозы нет>>.
  \item \textbf{№\,6.} $L_R = 1$ ($x=5$), слон $(1,8)$ не на
        вертикали $x=5 \Rightarrow A_R = 1$; $|5-1| = |4-8| = 4
        \Rightarrow L_B = 1$, ладья $(5,8)$ не лежит на диагонали
        $a8$--$e4$ ($|5-5| = 0 \neq |4-8| = 4$) $\Rightarrow Z_B = 0
        \Rightarrow A_B = 1$. Вывод: <<Угроза от ладьи и слона>>.
\end{itemize}

Программа при прогоне этих наборов выдала те же результаты.
Дополнительно проверены случаи некорректного ввода (координаты 0 и~9,
совпадение клеток) --- программа корректно завершается с сообщением об
ошибке. \textbf{Расхождений между математической моделью и программой
не обнаружено.}

% ------------------------------------------------------------
% 8. ЗАКЛЮЧЕНИЕ
% ------------------------------------------------------------
\section*{8. Заключение}

В ходе выполнения работы по пунктам 2--6:
\begin{itemize}
  \item по п.\,2: постановка задачи 2.25 зафиксирована дословно по
        сборнику А.\,Г.\,Юркина; уточнены форматы ввода и вывода;
  \item по п.\,3: построена вербальная модель --- последовательность
        из семи шагов, включая правило перекрытия линии атаки;
  \item по п.\,4: построена математическая модель --- формулы условий
        боя $L_R$, $L_B$, перекрытия $Z_R$, $Z_B$ и итоговых признаков
        $A_R$, $A_B$; оценена сложность $O(1)$;
  \item по п.\,5: разработана блок-схема, декомпозирующая задачу на
        ввод, две группы проверок и ступенчатый вывод результата;
  \item по п.\,6: реализована программа на C++, структура которой
        повторяет математическую модель (отдельная функция на каждый
        признак), что снижает вероятность ошибки реализации.
\end{itemize}
Поставленная задача решена полностью: программа определяет наличие
угрозы королю и устанавливает её источник.

% ------------------------------------------------------------
% 9. ВЫВОДЫ
% ------------------------------------------------------------
\section*{9. Выводы}

\begin{enumerate}
  \item Разработана, реализована и протестирована программа на языке
        C++, решающая задачу 2.25: определение угрозы белому королю
        со стороны чёрной ладьи и чёрного слона.
  \item Алгоритм относится к классу линейных: циклы отсутствуют,
        объём вычислений фиксирован; сложность $O(1)$ по времени и
        по памяти.
  \item Построены вербальная и математические модели, блок-схема;
        согласованность артефактов подтверждена шестью контрольными
        примерами, включая случаи перекрытия линии атаки.
  \item Программа выполняет контроль корректности ввода (границы
        доски, совпадение клеток), что делает её устойчивой к ошибкам
        пользователя.
  \item Математическая модель легко модифицируется: при отказе от
        правила перекрытия достаточно исключить слагаемые $Z_R$ и
        $Z_B$; при расширении набора фигур (ферзь, конь, пешки)
        добавляются новые признаки $A_i$.
  \item Получены навыки формализации неформальной задачи,
        декомпозиции алгоритма, связки <<формула --- функция>> и
        тестирования по математической модели.
\end{enumerate}

% ------------------------------------------------------------
% 10. СПИСОК ИСТОЧНИКОВ
% ------------------------------------------------------------
\section*{10. Список использованных источников}

\begin{enumerate}
  \item Юркин А.\,Г. Сборник задач по программированию : учебное
        пособие. --- Иркутск : ИГУ, 2017. --- Раздел 2, задача 2.25.
        URL: http://itam.irk.ru/biblioavia/catalog/1505188362_527.pdf
        (дата обращения: 2026 г.).
  \item Страуструп Б. Язык программирования C++. 4-е изд. --- СПб. :
        Питер, 2020. --- 1328 с.
  \item Керниган Б., Ритчи Д. Язык программирования Си. 2-е изд. ---
        СПб. : Питер, 2020. --- 208 с.
  \item cppreference.com --- справочник по стандартной библиотеке C++.
  \item Методические указания к выполнению лабораторных работ по
        дисциплине <<Программирование>>. Раздел <<Локальная разработка.
        Линейные алгоритмы>>.
\end{enumerate}

% ------------------------------------------------------------
% 11. ПРИЛОЖЕНИЕ А
% ------------------------------------------------------------
\newpage
\section*{11. Приложение А. Полная распечатка кода}

Файл \texttt{main.cpp}. Компиляция:
\texttt{g++ -std=c++17 -Wall -o chess\_threat main.cpp}

\begin{Verbatim}[frame=single]
#include <iostream>
#include <cstdlib>
#include <string>

using namespace std;

struct Cell {
    int x;
    int y;
};

bool onBoard(const Cell& c) {
    return c.x >= 1 && c.x <= 8 && c.y >= 1 && c.y <= 8;
}

bool sameCell(const Cell& a, const Cell& b) {
    return a.x == b.x && a.y == b.y;
}

bool onDiagonal(const Cell& a, const Cell& b) {
    return abs(a.x - b.x) == abs(a.y - b.y) && a.x != b.x;
}

bool betweenStraight(const Cell& p, const Cell& a, const Cell& b) {
    if (a.x == b.x && p.x == a.x)
        return (a.y < p.y && p.y < b.y) || (b.y < p.y && p.y < a.y);
    if (a.y == b.y && p.y == a.y)
        return (a.x < p.x && p.x < b.x) || (b.x < p.x && p.x < a.x);
    return false;
}

bool betweenDiagonal(const Cell& p, const Cell& a, const Cell& b) {
    return onDiagonal(a, p) && onDiagonal(p, b) &&
           abs(a.x - p.x) < abs(a.x - b.x) &&
           abs(a.y - p.y) < abs(a.y - b.y);
}

bool rookThreat(const Cell& k, const Cell& r, const Cell& b) {
    if (k.x != r.x && k.y != r.y) return false;
    if (betweenStraight(b, k, r)) return false;
    return true;
}

bool bishopThreat(const Cell& k, const Cell& b, const Cell& r) {
    if (!onDiagonal(k, b))        return false;
    if (betweenDiagonal(r, k, b)) return false;
    return true;
}

Cell readCell(const string& who) {
    Cell c;
    cout << "Введите координаты " << who << " (x y, 1..8): ";
    cin >> c.x >> c.y;
    if (!cin || !onBoard(c)) {
        cerr << "Ошибка: клетка вне доски." << endl;
        exit(1);
    }
    return c;
}

int main() {
    setlocale(LC_ALL, "Russian");

    Cell K = readCell("белого короля");
    Cell R = readCell("черной ладьи");
    Cell B = readCell("черного слона");

    if (sameCell(K, R) || sameCell(K, B) || sameCell(R, B)) {
        cerr << "Ошибка: фигуры не могут стоять на одной клетке." << endl;
        return 1;
    }

    bool rook   = rookThreat(K, R, B);
    bool bishop = bishopThreat(K, B, R);

    if (rook && bishop)  cout << "Угроза королю от ладьи и слона." << endl;
    else if (rook)       cout << "Угроза королю от ладьи." << endl;
    else if (bishop)     cout << "Угроза королю от слона." << endl;
    else                 cout << "Угрозы королю нет." << endl;

    return 0;
}
\end{Verbatim}

\end{document}
